\section{Extensión: Compensation Tickets}
\label{sec:extension}

\subsection{Extensión elegida y motivación}

De las cinco extensiones propuestas por el enunciado (Tabla~\ref{tab:opciones},
Sección~\ref{sec:analisis-extension}), el equipo eligió la
\textbf{Opción A: \textit{compensation tickets}}. La extensión exige
"modelar una tarea que cede antes de consumir el \textit{quantum} e inflar
temporalmente su valor según la fracción consumida", y responde a la
pregunta "¿recupera la participación de CPU a la que tenía derecho?". En el
paper original~\cite{Waldspurger1994}, este mecanismo compensa a un
servidor que se bloquea esperando E/S en nombre de un cliente y no debería
ser penalizado frente a clientes que nunca se bloquean: al reactivarse, el
servidor recibe temporalmente más boletos, proporcionales al inverso de la
fracción de su último \textit{quantum} que efectivamente usó, hasta agotar
el equivalente a un \textit{quantum} completo de CPU recibido en boletos
compensados.

Se eligió esta opción sobre las otras cuatro porque reutiliza el mismo
mecanismo cooperativo/quantum que el proyecto base de todas formas debía
implementar (no agrega una máquina de estados nueva ni una fuente de
eventos externa), porque su pregunta de evaluación se responde con el mismo
experimento de proporcionalidad que ya exige el proyecto (comparar
\texttt{observed\_share} con y sin compensación), y porque presenta menor
superficie de bugs de concurrencia nuevos que la Opción C (estados
cliente/servidor) o la D (doble contabilidad de monedas).

\subsection{El escenario debe forzarse}

Con el diseño base del planificador (modos cooperativo y quantum tal como
los especifica el enunciado), una tarea \emph{siempre} consume todo su
bloque asignado --- no existe ningún mecanismo para que ceda antes de
tiempo por cuenta propia. El escenario de cesión temprana que la extensión
necesita observar no puede surgir orgánicamente; hay que modelarlo
explícitamente, como pide el propio enunciado. Para eso se agregó un
parámetro opcional y exclusivo del experimento de la extensión:
\texttt{--yield-config <csv>}, un archivo con encabezado
\texttt{task\_id,yield\_percent} que le indica a tareas puntuales que cedan
tras correr solo una fracción de su bloque en vez de agotarlo. Ausente por
defecto, no tiene ningún efecto sobre los siete casos mínimos del
enunciado ni sobre el comportamiento base del planificador.

\subsection{Mecanismo}

El \textit{quantum} (o el bloque cooperativo) se mantiene fijo: lo que se
infla temporalmente son los \textbf{boletos} de la tarea que cedió temprano,
para darle mayor probabilidad de ganar los siguientes sorteos y así
compensar el turno reducido. Si una tarea consumió solo una fracción $f$ de
su bloque antes de ceder:
\begin{equation}
\text{tickets\_compensados} = \left\lceil \text{tickets\_base} \times \frac{1}{f} \right\rceil
\label{eq:compensacion}
\end{equation}
La compensación no es permanente: se mantiene mientras la tarea tenga una
\textbf{deuda} pendiente, definida como lo que le faltó para completar el
bloque original en la activación donde cedió. Cada vez que la tarea gana un
sorteo con boletos inflados, corre el bloque completo (ya no vuelve a ceder
mientras compensa) y descuenta lo corrido de la deuda; al llegar la deuda a
cero, \texttt{effective\_tickets} vuelve a \texttt{tickets\_base} y el ciclo
puede repetirse en la siguiente activación que gane.

\subsection{Decisiones de implementación}

Tres puntos que el diseño original dejaba abiertos a propósito se
resolvieron al implementar (detalle completo en
\texttt{decisiones\_diseno.md}, D15, del repositorio de conocimiento):

\begin{enumerate}
\item \textbf{La fracción $f$ se mide contra el bloque que decide el modo,
  no contra un quantum fijo.} Con dos modos de ejecución, $f$ se generaliza
  como $\text{unidades\_corridas} / \text{bloque\_decidido}$, donde
  \texttt{bloque\_decidido} es la misma cantidad que ya calculaba
  \texttt{decide\_slice\_units} para esa activación (quantum fijo, o el
  bloque cooperativo recortado por lo que falta) --- sin ramas de código por
  modo dentro de la extensión.
\item \textbf{El campo \texttt{tickets} de \texttt{Task} no se renombra.}
  Ya cumple el rol de "boletos base"; se agregó \texttt{effective\_tickets}
  como campo nuevo (igual a \texttt{tickets} salvo compensación activa) en
  vez de tocar los sitios que ya construían o leían \texttt{Task.tickets}
  desde el Milestone 1.
\item \textbf{La configuración de qué tareas ceden vive en un archivo
  aparte.} Nunca en el CSV principal de tareas: una bandera ausente por
  defecto es una garantía más fuerte de que los casos base no se ven
  afectados que una columna opcional que el parser tendría que aprender a
  ignorar.
\end{enumerate}

El control del experimento A/B (aislar el efecto de \emph{ceder} del efecto
de \emph{compensar}) tampoco estaba resuelto por el documento de diseño
original: se agregó \texttt{--disable-compensation}, que requiere
\texttt{--yield-config} y mantiene la cesión temprana activa pero nunca
infla \texttt{effective\_tickets} ---la tarea cede la misma fracción en cada
activación, indefinidamente, sin compensación--- de modo que la comparación
entre corridas con y sin esta bandera aísla exactamente el efecto de la
compensación en sí.

\subsection{Hipótesis}

\begin{itemize}
\item \textbf{H1.} Una tarea que cede voluntariamente antes de agotar su
  bloque, y que recibe \textit{compensation tickets} proporcionales a la
  fracción no consumida, alcanza un \texttt{observed\_share} significativamente
  más cercano a su share objetivo que la misma tarea bajo idénticas
  condiciones sin compensación.
\item \textbf{H0.} La compensación no tiene efecto medible sobre el
  \texttt{observed\_share} de una tarea que cede tempranamente.
\end{itemize}

El diseño experimental (control \texttt{--disable-compensation} vs.
tratamiento con compensación activa, misma semilla, mismo CSV, misma
ventana de despachos, repetido sobre múltiples semillas) y sus resultados
se presentan en la Sección~\ref{sec:resultados-extension}.
